% Метод Хука--Дживса: исследующий поиск по координатам + прыжок по образцу.
\begin{figure}[htbp]
\centering
\begin{tikzpicture}[scale=1.0, >={Stealth[length=2.2mm]}]

% линии уровня
\foreach \s in {1.0, 0.68, 0.36} {
  \draw[blue!55!black, rotate=-25] (0,0) ellipse ({2.6*\s} and {1.5*\s});
}
\fill[blue!55!black] (0,0) circle (1.2pt);
\node[font=\scriptsize, blue!55!black, anchor=west] at (0.1,0.02) {$x_*$};

% базовая точка и исследующий поиск
\coordinate (b0) at (1.55,-1.30);
\coordinate (e1) at (1.05,-1.30);   % шаг -h_1 e_1
\coordinate (b1) at (1.05,-0.90);   % шаг +h_2 e_2

\draw[orange!75!black, thick, ->] (b0) -- (e1);
\draw[orange!75!black, thick, ->] (e1) -- (b1);
\fill[orange!75!black] (b0) circle (1.8pt) (e1) circle (1.4pt)
                        (b1) circle (1.8pt);
\node[font=\scriptsize, anchor=west] at ($(b0)+(0.1,-0.05)$) {$x^{(0)}$};
\node[font=\scriptsize, anchor=north west] at ($(b1)+(0.08,-0.02)$) {$x^{(1)}$};
\node[font=\scriptsize, orange!60!black, anchor=north] at (1.3,-1.42)
  {исследующий поиск};

% прыжок по образцу: x^(2) = 2 x^(1) - x^(0)
\coordinate (p2) at (0.55,-0.50);
\draw[blue!45!black, very thick, ->] (b1) -- (p2);
\fill[blue!45!black] (p2) circle (1.8pt);
\node[font=\scriptsize, anchor=south east] at ($(p2)+(0.02,0.06)$)
  {$2x^{(1)} - x^{(0)}$};
\node[font=\scriptsize, blue!45!black, anchor=south west] at (0.62,-0.92)
  {поиск по образцу};

\end{tikzpicture}
\caption{Метод Хука--Дживса. Из базовой точки $x^{(0)}$ исследующий поиск
пробует сдвиг по каждой координате на $\pm h_i$ и принимает те, что уменьшают
$f$; получается новая базовая точка $x^{(1)}$. Затем делается прыжок по образцу
$x^{(2)} = 2x^{(1)} - x^{(0)}$ — повтор уже удачного смещения без пробных
вычислений. Если прыжок неудачен, метод откатывается к последней удачной точке и
повторяет исследующий поиск с уменьшенными $h_i$.}
\label{fig:hooke-jeeves}
\end{figure}
